\section{落地建议：团队如何构建抗噪声评测流水线}

\subsection{一个可执行的评测 SOP}
结合本讲内容，可以整理出一个可落地的评测 SOP：先定义任务层级和关键指标，再固定协议版本，然后做重复评测并报告置信区间，最后再做跨模板和跨环境鲁棒性验证。这个顺序能避免 ``先看排行榜再补解释'' 的被动流程。

\begin{table}[H]
\centering
\begin{tabular}{p{1.2cm}p{3.7cm}p{5.3cm}}
\toprule
\textbf{阶段} & \textbf{动作} & \textbf{交付物} \\
\midrule
1 & 任务定义与分桶 & task schema、难度分桶、目标指标 \\
2 & 协议冻结 & prompt 模板、decode 参数、judge 配置版本 \\
3 & 重复评测 & mean/std/CI、seed sweep 结果 \\
4 & 鲁棒性检查 & 模板扰动、环境扰动、judge 一致性报告 \\
5 & 上线判定 & 显著性门槛、回归报警规则、回滚策略 \\
\bottomrule
\end{tabular}
\caption{抗噪声 benchmark 流水线建议}
\end{table}

\subsection{当资源有限时，优先做什么}
课程里一个务实观点是：很多团队不是不知道统计方法，而是没有预算全做。若只能选几件事，优先级建议是：重复评测（k 次运行）> 协议版本固定 > 最小置信区间报告 > 难度分桶。这四项能用最小成本换来最大可解释性提升。

\begin{importantbox}{低成本高收益实践包}
\begin{itemize}
    \item 同一配置至少跑 5 次并报均值与标准差；
    \item 每次发布写明协议版本与 seed；
    \item 对最关键 20\% 样本做人评抽检；
    \item 对 hard split 单独汇报，不与总体平均混报。
\end{itemize}
\end{importantbox}

\subsection{本章小结}
抗噪声评测不是大厂特权，而是一套可分层实施的方法。即使资源有限，也可以通过重复评测、协议冻结和最小统计报告显著提高结论可信度。

